%% 0_abstract
\begin{abstract}
Radio map reconstruction recovers a dense signal-strength field from sparse
measurements, and progress on it has been pursued almost entirely by varying the
architecture. We show that architecture is not the binding constraint. Evaluating
nine learned reconstructors spanning convolutional, cascaded, state-space,
transformer, adversarial and diffusion families under a single protocol, we find
that they separate in line of sight, with a coefficient of variation of $20.1\%$, and barely
separate at all in non line of sight, at $9.6\%$, even though non line
of sight dominates the total error. A plateau this indifferent to design indicates a missing
input rather than insufficient capacity. The missing input is a line integral,
namely the fraction of the transmitter-to-pixel ray that lies inside building
material. Supplying it as one channel costs under $600$ parameters and pushes both error
modes below the entire panel. A controlled $2 \times 2$ shows that this geometry and
a deeper cascaded backbone are substitutes rather than complements, delivering
$2.87$~dB jointly where their separate gains predict $3.54$~dB, so a plain network
given the geometry beats a cascade denied it. The same result explains a transfer
failure. The integral is anchored at the transmitter, and on real indoor cellular
crowdsensing only $17.9\%$ of cells yield a locatable source. We show that
marginalising the anchor over measurement locations destroys the directional
structure that makes the channel work, and that on real measurements classical
interpolation still beats every learned reconstructor we trained. We validate the
underlying wall-attenuation premise directly on measurements against an open-plan
negative control.
\end{abstract}
